

\paragraph{Gauge transformation}
\label{p:sec:hs:sextic:gauge}

Under the gauge transformation~\eqref{p:eq:hs1:gauge-psi} the truncated action is invariant up to the order of the truncation, $\delta_\Lambda S_{\le 6} = O(\gs^5)$, which should be reproduced by the cubic action.
The two auxiliary fields enter $\caction[\le 6]$ through two completed squares, and we read off their transformation laws by demanding that each square be invariant, as at the quintic order.

The $3$-string field enters only through $\Delta S_6$, which is the square of the fluctuation:
\begin{equation}
  \label{p:eq:hs6:deviation-3}
  \sigma_3^{(4)}
    := \Sigma_3 - \Sigma_3^{(4)*}
    = \Sigma_3 - \Psi \otimes \Sigma_2,
  \qquad
  \Delta S_6
    = - \frac{\gs^4}{6!} \, \pair[3]{\sigma_3^{(4)}}{\wprod{3}{3}(\sigma_3^{(4)})},
\end{equation}
with the field-independent kinetic operator~\eqref{p:eq:hs6:kinetic}.
Its variation is:
\begin{equation}
  \label{p:eq:hs6:square-variation-3}
  \delta_\Lambda \Delta S_6
    = - \frac{2 \gs^4}{6!} \, \pair[3]{\delta_\Lambda \sigma_3^{(4)}}{\wprod{3}{3}(\sigma_3^{(4)})},
\end{equation}
with no term from the operator.
It vanishes for all $\sigma_3^{(4)}$ when $\delta_\Lambda \sigma_3^{(4)}$ lies in the kernel of $\wprod{3}{3}$, that is, when it vanishes on the quotient~\eqref{p:eq:hs6:quotient}.
We take this transformation law.
As for $\Sigma_2$ below, it is unique up to the trivial symmetries of \cref{p:sec:gauge}.

The solution~\eqref{p:eq:hs6:saddle-sigma3} and the propagator of $\Sigma_3$ rest on the invertibility of $\wprod{3}{3}$ on the quotient, assumed in~\eqref{p:eq:hs6:invertibility-3}.
We thus take the fluctuation of the $3$-string field inert, and $\Sigma_3$ transforms as its solution:
\begin{equation}
  \label{p:eq:hs6:gauge-sigma3}
  \delta_\Lambda \sigma_3^{(4)}
    = 0,
  \qquad
  \delta_\Lambda \Sigma_3
    = \delta_\Lambda \Psi \otimes \Sigma_2 + \Psi \otimes \delta_\Lambda \Sigma_2.
\end{equation}
There is no rotation term: the rotation at quintic order compensated the variation of the dressed operator, which depends on $\Psi$, and $\wprod{3}{3}$ does not depend on $\Psi$.
This is the quartic mechanism one level up, with $\Psi \otimes \Sigma_2$ in place of $\Psi^2$.

Once $\Sigma_3$ is eliminated, the action~\eqref{p:eq:hs6:eliminate-sigma3} is quadratic in $\Sigma_2$, and its quadratic part defines the dressed kinetic operator of $\Sigma_2$ at this order:
\begin{equation}
  \label{p:eq:hs6:dressed-2}
  \pair[2]{A}{\kin{2}{\Psi}{4}(B)}
    := \pair[2]{A}{\kin{2}{\Psi}{3}(B)}
      - \frac{\gs^2}{30} \, \dvertex{6}(\Psi \otimes A, \Psi \otimes B),
\end{equation}
for $A, B \in \Hn{2}$, so that the terms of $\caction[\le 6] |_{\Sigma_3 = \Sigma_3^{(4)*}}$ quadratic in $\Sigma_2$ read
\begin{equation}
  - \frac{\gs^2}{4!} \, \pair[2]{\Sigma_2}{\kin{2}{\Psi}{4}(\Sigma_2)},
\end{equation}
as in~\eqref{p:eq:hs5:F-dressed}.
The correction is the $\Sigma_3$ exchange between two of the cubic couplings $\dvertex{6}(\Psi \otimes \Sigma_2, \Sigma_3)$, collapsed to a contact term by the elimination of $\Sigma_3$.
The operator is symmetric with respect to the $2$-string pairing and depends on $\Psi$ through both terms.
The solution and the fluctuation are defined as at the quintic order:
\begin{equation}
  \label{p:eq:hs6:saddle-sigma2}
  \Sigma_2^{(4)*}
    := \big( \kin{2}{\Psi}{4} \big)^{-1} \circ \sprod{2}{2}(\Psi^2),
  \qquad
  \sigma_2^{(4)}
    := \Sigma_2 - \Sigma_2^{(4)*},
\end{equation}
the inverse being a Neumann series in $\gs$ around $\sprod{2}{2}^{-1}$ with one more kind of insertion, the order-$\gs^2$ term of~\eqref{p:eq:hs6:dressed-2}.
Completing the square as in~\eqref{p:eq:hs5:square}, we find that the cubic action is quadratic in $\sigma_2^{(4)}$ with the operator $\kin{2}{\Psi}{4}$.
We demand that this quadratic form be invariant.
Its variation has a term from the fluctuation and one from the $\Psi$ inside the operator, and it vanishes for:
\begin{equation}
  \label{p:eq:hs6:gauge-sigma2}
  \delta_\Lambda \sigma_2^{(4)}
    = - \frac{1}{2} \, \big( \kin{2}{\Psi}{4} \big)^{-1} \circ \big( \delta_\Lambda \kin{2}{\Psi}{4} \big) (\sigma_2^{(4)}),
  \qquad
  \delta_\Lambda \Sigma_2
    = \delta_\Lambda \Sigma_2^{(4)*} + \delta_\Lambda \sigma_2^{(4)},
\end{equation}
where $\delta_\Lambda \kin{2}{\Psi}{4}$ is the variation of the operator~\eqref{p:eq:hs6:dressed-2} through its argument $\Psi$, by the symmetry of the operator.
At order $\gs$ the variation of the operator is $\sprod{3}{2}(\delta_\Lambda \Psi, \cdot)$, and~\eqref{p:eq:hs6:gauge-sigma2} is the transformation law~\eqref{p:eq:hs5:gauge-deviation} of the quintic order.
The order-$\gs^2$ term is the variation of the sextic correction.
This is the transformation law of \cref{p:sec:gauge}, written there as~\eqref{p:eq:gauge:law-quintic} with the dressed operator of the order at hand.
It has the structure of the quintic one: the auxiliary field transforms as its solution, plus a rotation of the fluctuation which is linear in it, carries no parameter of its own and vanishes on-shell.
The transformation law is unique up to the trivial symmetries of \cref{p:sec:gauge}.

With the transformation laws~\eqref{p:eq:hs6:gauge-sigma3} and~\eqref{p:eq:hs6:gauge-sigma2}, both squares are exactly invariant.
The off-shell identity behind~\eqref{p:eq:hs6:eliminate}, the analogue of~\eqref{p:eq:hs5:identity}, reads:
\begin{equation}
\caction[\le 6]
  = S_{\le 6}
    + (\text{exchange terms of order } \gs^5)
    + (\text{squares}).
\end{equation}
It gives:
\begin{equation}
  \label{p:eq:hs6:delta-S-equal}
  \delta_\Lambda \caction[\le 6] - \delta_\Lambda S_{\le 6}
    = O(\gs^5),
\end{equation}
identically, off-shell, for any $\delta_\Lambda \Psi$: the cubic action is invariant to the order of the truncation it reproduces, as at the two previous orders.
The algebra of these transformations is that of closed string field theory through this order (\cref{p:sec:gauge}).


\paragraph{A second \texorpdfstring{$2$}{2}-string field}
\label{p:sec:hs:sextic:partner}

\subparagraph{Localizing the action.}

The reduced vertex is one way of subtracting the exchange term.
  Another way keeps every vertex elementary, at the price of a second $2$-string field.
  We give a field $\Upsilon_2 \in \Hn{2} / \ker \sprod{2}{2}$, with the ghost number and the parity of $\Sigma_2$, a quadratic term and a single elementary coupling, with coefficients $a$ and $b$:
  \begin{equation}
  \label{p:eq:hs6:partner}
  S_\Upsilon
    := \frac{\gs^2}{4!} \, a \, \vertex{4}(\Upsilon_2^2)
      + \frac{\gs^3}{5!} \, b \, \vertex{5}(\Upsilon_2, \Sigma_2, \Psi),
  \qquad
  \frac{b^2}{a}
    = 4.
\end{equation}
  Since $\Upsilon_2$ has no source at order $\gs^2$, its classical value starts at order $\gs$.
  Its own exchange, two quintic vertices joined by an $\Upsilon_2$ line, is $- b^2 / (4 a)$ times the exchange term of $\Sigma_2$.
  The two cancel for $b^2 / a = 4$, and we take $(a, b) = (1, 2)$, the $1 : 2$ ratio of the quadratic and source terms of $\Sigma_2$ in~\eqref{p:eq:hs4:S4tilde}, so that the quadratic form of $\Upsilon_2$ has the sign opposite to that of $\Sigma_2$.
  The action
  \begin{equation}
  \caction[\le 5] + S_\Upsilon + \frac{\gs^4}{6!} \, \vertex{6}(\Sigma_2^3)
\end{equation}
  then reproduces $S_{\le 6}$ through order $\gs^4$ with the fields $(\Psi, \Sigma_2, \Upsilon_2)$, elementary vertices only, and no $\Sigma_3$.
  This holds at tree level, for generic off-shell $\Psi$, under the same assumptions as the construction with $\Sigma_3$, and it was checked in the finite-dimensional models of \cref{p:app:finite-models} with generic vertices.
  For a real field the opposite sign of the quadratic form is required, and the sign is harmless at tree level.
  We leave open whether the sign has a physical consequence beyond tree level, $\Upsilon_2$ being an auxiliary field without propagating states, as well as the gauge transformation of $\Upsilon_2$ and the completed-square form of its sector.
  We give the construction in \cref{p:app:wrong-sign}.

A same-sign field $\Xi_2$ can instead absorb the inverse of $\sprod{2}{2}$ in the kinetic term of $\Sigma_3$, as the Liouville field makes the induced gravitational action local.
It cannot do the same in the coupling of $\Sigma_3$ to $(\Psi, \Sigma_2)$ at the same time: one field cannot serve two terms without producing a cross term.
A fully local action would need several of them, as in the intermediate-field decompositions of~\cite{Rivasseau:2010:LoopVertexExpansion,Lionni:2018:NoteIntermediateField}.
This too is in \cref{p:app:wrong-sign}.
